% Section 3 of the manuscript: what can be logged without reading text.
% This analysis is not part of the registration. Every statement rests on a document listed in the supplement
% (with its URL, an archived copy, its date and the date it was retrieved) or is marked as an observation.

\section{What can be logged without reading text}
\label{sec:observability}

A program that logs the conversion of another application without reading text needs an interface through which
it can learn, from outside that application, when a conversion begins and ends. This section examines which of the
documented public interfaces of Windows and macOS offer this. It is an analysis of documents, not part of the
registration, and its claims are stated for the interfaces examined.

\subsection{Method}
\label{sec:observability-method}

The interfaces were chosen from the sections of each system's developer documentation on text input, input methods,
accessibility, event hooks and windows, and from the interfaces that earlier loggers use (the low-level keyboard
hooks of keystroke loggers and the plug-ins inside particular applications of the hybrid logger;
Section~\ref{sec:related}). On Windows the interfaces examined were the Text Services Framework, the input-method
manager, message hooks and the low-level keyboard hook, accessibility events and the accessible objects of the input
method's windows, UI Automation, and the desktop compositor. On macOS they were the text-input client protocol, the
accessibility attributes and notifications of both the C interface and the AppKit framework, the Input Method Kit,
the window list, the text input sources and event taps.

Each document was read in the version published on the developer site on 5 October 2026, when it was retrieved and
an archived copy of it was made; the list of documents in the supplement gives, for each, its URL, the URL of the
archived copy, the date of its last update as the site states it (the Apple documentation states none, but gives the
system version in which each item was introduced) and the date of retrieval. The documents describe the current
systems; the measurement used Windows 11 25H2 and macOS 27.0.1. For each interface the
criterion was whether a program that runs outside the target application can learn through it when a conversion is
open, without receiving the composition, the candidates or any other text. An interface that requires code to be
loaded into the target application does not meet the criterion, because such code runs inside the application and
can read its text. A statement that rests on observing a system rather than on a document is marked as an
observation, with the system and its version as far as they were recorded.

\subsection{Windows}
\label{sec:observability-windows}

The Text Services Framework is the interface through which the Microsoft input method delivers its composition to
an application. Its thread manager is activated for the calling thread \citep{ms-tsf-activate}, and its sinks are
installed on a source obtained from the thread manager or from a context of the calling process
\citep{ms-tsf-advisesink}; no parameter of these calls names another process. The documentation for assistive
technology states that ``to access the composition string managed in TSF, the assistive technology products must
create the ITfThreadMgr in the process of creating the target TSF-enabled applications''
\citep{ms-ime-tsf-accessibility}. The text services of the framework, which provide text to applications, are
implemented as in-process servers registered with the framework \citep{ms-tsf-architecture}, so code that takes this
route runs inside the application. The input-method manager assigns an input context to each thread by default, and
an application retrieves the context associated with a specified window \citep{ms-imm-inputcontext}; its
composition string is the text being converted \citep{ms-imm-getcompositionstring}.
Message hooks that would intercept the input-method messages of other windows load a library into the target
processes \citep{ms-setwindowshookex}.

Two documented routes remain that do not load code into the target and do not carry text. First, accessibility
events report that a window of the input method has become visible, has changed its size or position, or has become
hidden \citep{ms-eventconstants}. A program can receive these events out of context, so that its hook
function is not mapped into the address space of the application that raised them \citep{ms-outofcontext,
ms-setwineventhook}, and the arguments of the callback are the handle of the hook, the event, a window handle, two
object identifiers, a thread identifier and a time \citep{ms-wineventproc}: none of them is a string. The
documentation for assistive technology, written for Windows XP, describes the other direction: following the events
of the candidate window to its accessible objects reads the candidate strings, and the composition string can be read
in the same way where the input method draws it \citep{ms-ime-tsf-accessibility}. Second, accessibility events
report that a window has been cloaked or uncloaked \citep{ms-eventconstants}, and the desktop compositor tells whether a window is cloaked
\citep{ms-dwm-attribute}. During the development of the logger, on a Windows~11 machine whose build was not recorded,
the candidate window of the Microsoft input method was shown and hidden by uncloaking and cloaking it, and the window
events of the input method did not arrive (an observation made in August 2026). The logger follows both routes
(Section~\ref{sec:logger}).

\subsection{macOS}
\label{sec:observability-macos}

On macOS the marked text, the reading being converted, is held by the text view of the target application, which
implements the text-input client protocol to receive it \citep{apple-nstextinputclient}; the documentation examined
describes no call by which a program outside the application can query that protocol. An input method is a program of its own that creates an input
controller for each input session of a client application \citep{apple-imkinputcontroller}, and the Input Method Kit
is integrated with the system's text services \citep{apple-imk}; the documentation examined describes no part of an
input method that runs inside the client application. Assistive applications reach other applications through the accessibility interface
\citep{apple-axuielement}. Of its C constants listed in the documentation (101 attributes, 9 parameterized attributes
and 25 notifications), none concerns marked text or a composition
\citep{apple-ax-attributes,apple-ax-parameterized,apple-ax-notifications}; the header files of the software
development kit of macOS 27.0 define more of them (138 attributes, 14 parameterized attributes and 37 notifications,
counted by the ending of their names), and none of these concerns marked text either \citep{apple-sdk-27}. The
constants of the AppKit framework (131 attributes, 18 parameterized attributes and 42 notifications, the same counts
in the documentation and in the header file of the kit) include three that do: an attribute for the marked
range of a text input and two notifications for the beginning and end of a marking session, available since macOS
10.6 and listed without a description
\citep{apple-nsaccessibility-attribute,apple-nsaccessibility-parameterized,apple-nsaccessibility-notification}. They
are taken up below.

The list of windows is returned as it stands when it is requested, with information about each window such as its
bounds and its identifier \citep{apple-cgwindowlist}, so a program can follow it only by polling. It lists windows, not marked text: the reading
that is typed before conversion is shown by the application's text view, and a window of the input method exists only
when the input method shows one. The current keyboard input source \citep{apple-textinputsources} identifies the
selected input source and input mode, and none of its properties describes a composition; in the header file of
the kit of macOS 27.0 none of its functions or constants concerns a composition either \citep{apple-sdk-27}. An event tap receives the
key-down and key-up events of the session \citep{apple-cgevent-tapcreate}. The logger therefore infers sessions from
the key presses and the input source (Section~\ref{sec:logger}).

\subsection{Routes that depend on the application}
\label{sec:observability-application}

On both systems an application's own accessibility implementation can report a conversion without text. On Windows
the text-edit pattern of UI Automation defines an event for a change of the conversion target, which the elements that
support the pattern are required to raise \citep{ms-uia-textedit}, and the corresponding accessibility event carries
no string \citep{ms-eventconstants,ms-wineventproc}; the same pattern's event for a finished composition, however, carries the
finalized string \citep{ms-uia-textedit}. On macOS the attribute and the two notifications named above report the
range of the marked text (a location and a length) and the beginning and end of a marking session; the source code of
the WebKit engine posts the two notifications when a composition begins and ends and reports the marked range as a
location and a length \citep{webkit2026}. Whether these items are provided depends on the application and its
toolkit: on Windows the events are raised by the elements of an application that support the pattern
\citep{ms-uia-textedit}, and the macOS documentation does not state which elements provide the three items. A
logger meant to work in any application cannot rely on them, and the logger does not use them. Whether they report
the sessions of the reference input fields, or of the applications that writing research uses, was not examined.

The routes that do not depend on the application therefore differ between the systems. On Windows they are the window
events of the input method and the cloaking of its windows; on macOS none of the documented interfaces examined
reports a conversion, and the logger infers it from key presses. The registration states more broadly that no public
interface of macOS reports the conversion state of another application; the documents support the narrower
statement made here. Whether a remaining route works depends on the input method and the application as well: the
cloaking path follows windows of the input method's class, and an application that draws the candidates itself may
show no window of the input method (Section~\ref{sec:limitations}).

\subsection{Consequence}
\label{sec:observability-consequence}

The two systems log different constructs: on Windows a period during which a window of the input method is visible,
and on macOS a segment of kana input inferred from the keys. Their measures are therefore defined differently, they
are validated separately against references of their own system (Section~\ref{sec:method-reference}), and they are not
compared as equal. Table~\ref{tab:interfaces} summarises the interfaces examined.

% A landscape long table (main.tex): a sideways float of this length runs off the rotated page.
\begin{landscape}
{\tablebodyfont
\begin{longtable}{L{0.26}L{0.35}L{0.08}L{0.18}L{0.11}}
\caption{Interfaces examined for logging the conversion of another application. ``Text'' is whether the interface
delivers the composition, the candidates or other text to the program that uses it}
\label{tab:interfaces}\\
\toprule
Interface & What it exposes to a program outside the application & Text & Used by the logger & Source \\
\midrule
\endfirsthead
\tablecontinued{5}
\toprule
Interface & What it exposes to a program outside the application & Text & Used by the logger & Source \\
\midrule
\endhead
\multicolumn{5}{l}{\emph{Windows}} \\*
Text Services Framework (thread manager and sinks) & The composition of the calling thread; for another application only from inside its process & Yes & No & \citet{ms-tsf-activate,ms-tsf-advisesink,ms-ime-tsf-accessibility} \\
Text service & Text of the application in which it runs, as an in-process server & Yes & No & \citet{ms-tsf-architecture} \\
Input-method manager & The input context of a thread and its composition string & Yes & No & \citet{ms-imm-inputcontext,ms-imm-getcompositionstring} \\
Message hooks & Messages of other windows, through a library loaded into their processes & Yes & No & \citet{ms-setwindowshookex} \\
Low-level keyboard hook & The key code and time of every key event & No & Key rows; \texttt{nav} and \texttt{abandon} & \citet{ms-setwindowshookex} \\
Input-method window events, out of context & Display, change and dismissal of a window of the input method, system-wide & No & Sessions (first path) & \citet{ms-eventconstants,ms-outofcontext,ms-wineventproc} \\
Cloak events and the compositor's cloak state & Whether a window is cloaked; the window's class & No & Sessions (second path) & \citet{ms-eventconstants,ms-dwm-attribute} \\
Accessible objects of the input method's windows & Candidate strings; the composition string where the input method draws it & Yes & No & \citet{ms-ime-tsf-accessibility} \\
UI Automation text-edit pattern, finished composition & The finalized string & Yes & No & \citet{ms-uia-textedit} \\
UI Automation text-edit pattern, conversion target & A change of the conversion target, if the application raises it & No & No & \citet{ms-uia-textedit,ms-eventconstants} \\
Focused element and top-level windows & Properties of the focused element and of the top-level windows; the logger reads class names only & No, for class names & Secondary signals of a candidate window & \citet{ms-uia-getfocusedelement,ms-uia-classname,ms-enumwindows} \\
\midrule
\multicolumn{5}{l}{\emph{macOS}} \\*
Text-input client protocol & Marked and inserted text, inside the application only & Yes & No & \citet{apple-nstextinputclient} \\
Input Method Kit & The input method's own sessions with its clients & Yes & No & \citet{apple-imk,apple-imkinputcontroller} \\
Accessibility, C constants & No item on marked text & No & No & \citet{apple-ax-attributes,apple-ax-parameterized,apple-ax-notifications} \\
Accessibility, AppKit marked range and marking session & The range of the marked text and the beginning and end of a marking session, if the application provides them & No & No & \citet{apple-nsaccessibility-attribute,apple-nsaccessibility-notification,webkit2026} \\
Window list & Information about each window, such as its bounds, when polled & No & No & \citet{apple-cgwindowlist} \\
Text input sources & The selected input source and input mode & No & Whether a Japanese input source is selected & \citet{apple-textinputsources} \\
Event tap & The key code and time of every key event & No & Key rows; sessions by inference & \citet{apple-cgevent-tapcreate} \\
\bottomrule
\end{longtable}
}
\end{landscape}
